\documentclass[11pt]{article}
\usepackage[margin=1in]{geometry}
\usepackage{amsmath,amssymb}
\usepackage{booktabs}
\usepackage{array}
\usepackage{enumitem}
\usepackage{xcolor}
\usepackage{hyperref}
\hypersetup{colorlinks=true,linkcolor=blue,citecolor=blue}
\setlength{\parskip}{3pt}
\emergencystretch=1.5em

\title{The v3 Controller/Estimator Campaign on the Mecanum Digital-Twin Plant\\
\large What was run, why, and how to read it --- explainer companion to
\texttt{RESULTS\_v3.md}}
\author{Mecanum--PINN dynamics notes\\
\small KUKA youBot four-Mecanum platform, IMECE 2026 track\\
\small code: \texttt{hybrid\_ctrl\_v2/} (Julia closed-loop stack)}
\date{August 20, 2026}

\begin{document}
\maketitle

\begin{abstract}
\noindent
This note explains the v3 controller/estimator simulation campaign end to end: the
question it answers (\S1), the simulated plant, sensors and estimator (\S2), the
feasibility screening that defines which trajectories are fair to ask for (\S3), the
scalar objective every run minimizes (\S4), the five-stage protocol (\S5), the headline
findings (\S6), and the reading rules that keep the numbers honest (\S7). It is the
narrative companion to \texttt{hybrid\_ctrl\_v2/RESULTS\_v3.md}, which remains the
document of record for the numbers themselves.
\end{abstract}

\section{The question}
The project needs a defensible answer to: \emph{which control law should the
omnidirectional platform fly, and how much of its performance survives real state
estimation?} Three candidate laws are compared, chosen so the comparison is a two-step
ablation:

\begin{center}\small
\begin{tabular}{@{}l p{8.6cm} p{4.9cm} @{}}
\toprule
Controller & Structure & What it isolates\\
\midrule
\textbf{PID-FB} & cascaded IMC-reparameterized PID, 3 tuned $\lambda_{inner}$ per axis,
                  rest derived & industrial baseline\\
\textbf{PID-CT} & same cascade $+$ computed-torque feedforward (matched to ASMC's
                  equivalent control) & FB$\to$CT: value of the \emph{model}\\
\textbf{ASMC v2} & adaptive-gain supertwisting SMC, friction-circle gain ceiling
                  recomputed per tick & CT$\to$ASMC: value of \emph{adaptation}\\
\bottomrule
\end{tabular}
\end{center}

All three are tuned by the same black-box optimizer against the same scalar objective on
the same trajectory tiers, then evaluated through the same frozen state estimator.
Anything that makes one controller's objective differ from another's is treated as a bug,
not a detail. MPC was part of the original plan and is deferred out of this iteration
(diverges on four velocity-demanding trajectories; its \texttt{use\_ltv=false}-is-better
anomaly remains unresolved).

\section{The system under test}
\paragraph{Plant.} The 30-state voltage-input plant (\texttt{hybrid\_ctrl/plant.jl}):
slip--spin--LuGre contact physics (per-wheel bristle states, passive roller spin) driven
through a DC-motor model with back-EMF,
\begin{equation}
\tau = G\eta\,(K_t i - \tau_f),\qquad i=(V-K_bG\omega)/R_a,\qquad |V|\le24\ \text{V},
\end{equation}
with a 200 V/s slew clamp in the mixer. Actuator authority falls linearly with wheel
speed and vanishes at the no-load speed $V_{max}/(K_bG)=27.73$ rad/s --- the hard
feasibility wall of \S3. (Derived in
\texttt{docs/Mecanum\_Analytical\_Limits\_AxisVel\_AccelEnvelope\_BackEMF.tex}.)

\paragraph{Sensors (\texttt{:realistic} suite).} IMU and wheel encoders at 1000 Hz;
optical flow at 100 Hz with 2\% dropout; 100 Hz pose fixes with per-run biases and 0.5\%
outliers --- industrial white noise plus deliberately uncalibrated consumer-grade biases.
Neither the estimator nor (in the deployable configuration) the controllers ever see the
true state.

\paragraph{Estimator.} \texttt{ESKFEstimatorV3}: a 13-state error-state Kalman filter
with a yaw-acceleration random-walk state (a type-2 tracker, added to remove a
structural yaw-rate lag). It is tuned once and \textbf{frozen} before any controller
comparison --- the \emph{estimator-first methodology} --- so a control-law difference can
never be an observer difference in disguise.

\paragraph{Three feedback regimes --- not two.} This distinction runs through the whole
campaign; getting it wrong inverts conclusions:
\begin{enumerate}[nosep]
\item \textbf{clean oracle} --- perfect state feedback; an upper bound.
\item \textbf{noisy oracle} --- true state \emph{plus raw, unfiltered sensor noise}
      written into the controller's state estimate. There is \emph{no estimator} in this
      path; it is a pessimistic no-filtering bound, not a deployment model.
\item \textbf{ESKF} --- the deployable path: \texttt{:realistic} sensors through the
      frozen filter.
\end{enumerate}
The ESKF beats the noisy oracle \emph{by construction} (a tuned filter beats raw noise),
so ``ESKF vs noisy oracle'' is \textbf{not} the cost of state estimation --- that cost is
ESKF vs \emph{clean} oracle.

\section{Feasibility first: the v3 trajectory tiers}
The predecessor tier (\texttt{train12}) contained trajectories the hardware cannot fly,
and the tuner was spending most of its budget on them: \texttt{spiral\_orbit\_stress}
alone carried 72--81\% of the training tracking term while sitting at 85.4\% of the
motor's no-load speed. The v3 tiers (\texttt{train14\_v3}, \texttt{test\_v3}) screen
every candidate reference over its own full duration against two walls:
\begin{enumerate}[nosep]
\item \textbf{Motor / back-EMF (hard):} max per-wheel speed $(V_x\pm V_y\pm
(l{+}h)\dot\psi)/R$ against the 27.73 rad/s no-load speed; past $\sim$85\% of no-load no
gain law can arrest error growth. v3 targets $\le70\%$ for margin.
\item \textbf{Friction circle (soft):} per-wheel contact utilisation. Brief excursions
are momentary slip the LuGre plant absorbs --- measured to be harmless --- so this wall
informs but does not veto.
\end{enumerate}
The point of the rebuild: a controller comparison is only meaningful where the hardware
can actually go.

\section{The objective --- one scalar, three terms}
\begin{equation}
\text{score} \;=\; \text{tracking}_{v3} \;+\; 0.05\,(\text{ce}/V_{MAX}) \;+\;
0.36\,(\text{chatter}/\text{CHATTER}_{REF}).
\end{equation}
\textbf{tracking$_{v3}$}: six tolerance-normalised terms (position and heading $\times$
terminal / peak / typical), divided by a per-trajectory extent scaler $k_{traj}$. The v2
metric scored four samples out of 12k--60k per run (terminal and argmax only) and was
blind to mid-path oscillation; the time-normalised integral (``typical'') term is the
fix, with tolerances calibrated from measured achieved ratios. A term of 1.0 means
``exactly at target''. \textbf{ce}: mean summed wheel-voltage magnitude (energy proxy).
\textbf{chatter}: mean per-tick total variation of the wheel-voltage command against
$\text{CHATTER}_{REF}=0.8$ V/ms, the mixer's own slew clamp. This term exists because an
entire arc of v2 tuning kept railing on chatter the objective could not price (chatter
had been normalised by $V_{MAX}$ --- a voltage dividing a slew rate, $\sim$30--100$\times$
too weak). $\lambda_{chatter}=0.36$ is \emph{balance-preserving, not re-derived}: it
reproduces the accepted v2 tracking share (29.9\%) under the rescaled v3 metric, and the
objective code refuses $\lambda_{chatter}>1.44$ under \texttt{--metric v3} so a
carried-over v2 value cannot pass silently.

\paragraph{The estimator has its own objective.} A \emph{replay} score (estimator vs
ground truth on logged trajectories, no controller in the loop), same
tolerance-normalised shape, with velocity error \emph{during slip} priced twice
(\texttt{lam\_slip}=1.0): the point of an estimator on this platform is that its model
mismatch \emph{is} slip. It converged to $\sim$71\% a velocity-estimation objective,
matching what the controllers' inner loops consume.

\section{The protocol --- five stages, in order}
\begin{enumerate}[nosep]
\item \textbf{Estimator tuning (first, then frozen).} 5 seeds over a 7-dim log-scaled
      gain space, replay objective, \texttt{:realistic} sensors. All seeds converge to
      the same estimator within 0.23\% on a common noise realisation; seed 4 adopted and
      frozen for the whole campaign --- including tiers it was not tuned on.
\item \textbf{Feasibility screening} of trajectory tiers (\S3).
\item \textbf{Clean controller tuning} --- 5 seeds each, cold start, oracle-clean
      feedback, \texttt{train14\_v3}, convergence by plateau detection (the v1 tuner was
      caught terminating mid-descent, which invalidated its scatter claims).
\item \textbf{Noisy controller tuning} --- 4 seeds each, warm-started from the same run's
      clean optimum, feedback = noisy oracle. Run for all three controllers or none, so
      a noisy-tuned ASMC is never quoted against clean-tuned PIDs.
\item \textbf{Evaluation under three feedback regimes} --- the six converged configs
      (3 controllers $\times$ clean/noisy-tuned) on fresh sensor seeds 101--105:
      oracle-clean, oracle-noisy, and through the frozen ESKF (420 runs: $6\times14$
      trajectories $\times$ 5 seeds, zero non-converged). The held-out
      \texttt{test\_v3} leg (generalisation of the final estimator$+$controller stack)
      completes the campaign: 288 oracle sims $+$ 240 ESKF runs --- \S6 and
      \texttt{RESULTS\_v3.md} \S9.
\end{enumerate}
Throughout: one OS process per seed (the pipeline has global mutable state), BLAS pinned
to one thread for determinism, \texttt{stop\_reason} recorded per run, and outputs under
\texttt{runs\_*/} never hand-edited.

\section{Headline results}
Full tables and caveats in \texttt{RESULTS\_v3.md}; the structural findings:

\paragraph{PID-CT wins every regime.} Clean oracle 0.14340 vs ASMC 0.15318 (no seed
overlap; ordering PID-CT $<$ PID-FB $<$ ASMC); noisy oracle; and frozen-ESKF closed loop
0.34127, paired 69/70 and 70/70 against ASMC and PID-FB.

\paragraph{Noisy tuning does not transfer to deployment.} It helps against the noisy
oracle it was tuned on (ASMC 5/5, CT 5/5 paired) but the advantage vanishes through the
real ESKF, and it \emph{damaged} ASMC's heading $\sim$2$\times$ everywhere. Deploy
clean-tuned configs.

\paragraph{Real state estimation costs $+0.20$ of score, uniformly.} Against clean
oracle, essentially identical across all six configs (42.9--45.1\% of the
clean$\to$unfiltered gap recovered): the estimator's contribution is
controller-independent --- exactly what makes a frozen observer a fair basis for
comparison.

\paragraph{The error budget is a three-stage story.} Sensor injection (identical across
configs by construction) $\to$ ESKF survival (74--79\% of injected velocity error
removed; pose removal flat at $\sim$67.6\% on every trajectory) $\to$ controller
tracking, where essentially \emph{all} inter-controller spread lives. The estimator is
common mode; only the tracking stage discriminates control laws.

\paragraph{CT's one weakness is heading under sustained yaw.} Best heading on 9 of 14
trajectories ($\sim$1.1 mrad), collapsing to 29.8/37.6 mrad on exactly the two
sustained-yaw stress trajectories --- the same two where its yaw integral clamp fires on
17--20\% of ticks (0.0\% elsewhere). Three independent measurements agree; whether the
clamp binding \emph{harms} (vs merely binds) is not yet established.

\paragraph{Much of the tier is at the estimator's noise floor.} On 7 of 14 trajectories,
tracking error ($\sim$2--3 mm) sits at or below the estimator's own pose accuracy
($\sim$3.2 mm): further controller improvement there is unmeasurable with this sensor
suite, so the comparison is decided by the handful of trajectories where tracking rises
clear of the floor.

\paragraph{Heading is the ESKF's best channel; yaw rate is its worst.} Heading sees
6.9$\times$ error removal once an 11 ms startup transient on
non-zero-initial-heading references is accounted for; yaw rate sees only $\sim$52\%
removal --- the gyro is its only source, with no absolute fix to lean on.

\paragraph{Generalisation (\texttt{test\_v3}).} The held-out tier answers whether any of
the above transfers to trajectories the tuning never saw. Three cautions and one
headline:
\begin{itemize}[nosep]
\item \textbf{Anchor caveat.} \texttt{coupled\_vomega\_anchor} is combo-identical to the
      training tier's \texttt{coupled\_vomega\_stress} --- a deliberate cross-tier
      consistency check, which passes bit-exactly ($\pm0.000000$) --- and must be
      excluded from every generalisation claim. It is the worst-tracked entry in either
      tier, so leaving it in would dominate any mean. The tier therefore yields
      \textbf{7 genuinely held-out trajectories}, 3 on profiles never seen in tuning.
\item \textbf{PID-CT generalises as the winner, unambiguously.} Clean-tuned PID-CT beats
      ASMC and PID-FB \textbf{35/35 paired} each on the held-out set and is first on
      every individual held-out trajectory --- the campaign's most robust result.
\item \textbf{The ASMC-vs-PID-FB ordering does not generalise.} Their training-tier gap
      was carried by a single trajectory (\texttt{coupled\_vomega\_stress}, 66\% of
      ASMC's gap); excluding it, they are statistically indistinguishable (23/35 and
      8/35). The defensible claim: \textbf{ASMC and PID-FB are equivalent, and PID-CT
      beats both.}
\item \textbf{The estimator is tier-invariant.} Every error-budget channel lands within
      $\sim$3 percentage points of its training value, and on all 7 held-out
      trajectories every controller tracks to 1.6--4.2 mm / 1.0--2.5 mrad. The heading
      collapses occur only on the anchor --- a \emph{training} trajectory --- consistent
      with CT's yaw-clamp weakness being a sustained-yaw property, not a generalisation
      failure.
\end{itemize}

\section{Reading rules}
\begin{enumerate}[nosep]
\item \textbf{Never mix tuning-stage scores with eval scores} --- different noise
      realisations, overlapping ranges.
\item \textbf{Pair, don't average, across realisations} --- realisation variance
      ($\sim$6--7\% eval, $\sim$16\% tuning) swamps controller differences
      ($\sim$2--4\%). Every ordering claim is a paired 5/5 or 4/4, never a difference of
      marginal means.
\item \textbf{Best-seed selection under noise is confounded} --- seed 4 is both the best
      seed and the easiest realisation for all three controllers; applied identically,
      so cross-controller comparison holds.
\item \textbf{Absolute values, never \% degradation} --- a better clean tracker reads a
      worse \% for the same absolute result.
\item \textbf{Three feedback regimes, not interchangeable} --- \S2.
\end{enumerate}

\section{Reproduction map}
\begin{center}\footnotesize
\begin{tabular}{ll}
\toprule
Piece & Where\\
\midrule
plant (untouchable v1) & \texttt{hybrid\_ctrl/plant.jl, scheduler.jl, mixer.jl}\\
v2 controllers + search spaces & \texttt{hybrid\_ctrl\_v2/controllers\_v2.jl, tune\_controller\_v2.jl}\\
objective + trajectory tiers & \texttt{controller\_tuning/stage\_objective.jl, trajsets.jl}\\
tuning drivers & \texttt{run\_stage\_asmc\_v3.bat, run\_stage\_pid\_v3.bat} (noisy warm-starts)\\
estimator tuning & \texttt{estimator\_tuning/} (objective\_v2, param\_space\_v4, replay)\\
frozen estimator config & \texttt{runs\_estimator\_v4\_mu0p5\_train12/seed4/best\_config.json}\\
converged controller configs & \texttt{runs\_asmc\_v3/, runs\_pid\_v3/}\\
closed-loop ESKF eval & \texttt{eval\_v3\_eskf.jl}, data \texttt{runs\_eskf\_v3\_train14/}\\
diagnostics & \texttt{diag\_v3\_pid\_imax\_binding.jl, diag\_v3\_error\_budget.jl}\\
provenance / verification & \texttt{metric\_v3\_\{calib,verify\}.jl, scalers\_v3.jl, lambda\_v3\_verify.jl}\\
feasibility theory & \texttt{docs/Mecanum\_Analytical\_Limits\_AxisVel\_AccelEnvelope\_BackEMF.tex}\\
\bottomrule
\end{tabular}
\end{center}

\section{Deliberately not done (yet)}
\begin{itemize}[nosep]
\item \textbf{MPC} --- deferred (\S1).
\item \textbf{\texttt{test\_v3} generalisation} --- \textbf{done} (\S6): PID-CT
      generalises 35/35; ASMC/PID-FB tie. Caveat: the tier yields 7 held-out
      trajectories, not 8 (anchor is combo-identical to a training entry).
\item \textbf{CT yaw integral budget sensitivity} --- the clamp is measured to
      \emph{bind}; that it \emph{harms} is unmeasured.
\item \textbf{ESKF yaw-rate re-weighting} --- interacts with the CT heading finding; the
      obvious next experiment.
\item \textbf{$\hat\psi$ initialisation from the first pose fix} --- removes the 11 ms
      startup transient that makes heading statistics incomparable across
      mixed-initial-heading tiers.
\end{itemize}

\end{document}
